\chapter{Orbits and Rotation}
\label{chap:orbits-and-rotation}

The first two chapters tracked motion along a line. This one tracks motion that closes on itself --- an
orbit, a rotation --- where the record becomes a loop. The loop is the instrument's most important new
capability, the one every later proof rests on: a closed path it can carry a tally around, and from which the
tally can return \emph{changed}. The closing is the world's --- a motion that comes back to where it began does
so whether or not we are watching --- but the change the tally shows on return is measured against a frame we
hold fixed, and that frame is ours. That returning-changed is holonomy, the seed of every proof of a charged
loop to come: \emph{that} a loop carries a residue is the world's; \emph{which} way that residue falls --- which
sign, which sense --- is set by the frame we carry it against. The chapter introduces it gently --- first as
the closed cycle of an orbit, where the smooth-shadow ceiling debuts; then as the conserved tally of a
rotation, the case where the tally comes back unchanged; and finally, in Foucault's pendulum, as a residue a
loop genuinely carries.

\section{The Kepler Effect: an orbit is a closed cycle}
\label{se:kepler}

An orbit is the first motion the instrument reads as a loop. Where a line is traversed once, an orbit is a
closed admissible cycle --- a recorded motion that returns to where it began and is required to do so again,
period after period,
\begin{equation}
  \gamma(t + T) = \gamma(t).
\end{equation}
The closing is the world's: that the motion comes back, and keeps coming back, is enforced by nothing we do.
What is ours is the chart we lay it on --- the coordinates, the origin, the clock that stamps each return. It
is not equilibrium --- not a fixed point a body settles into --- but a stable failure to terminate, a loop of
length at least two that the record retraces forever. Because the path closes, the natural object for the
instrument is no longer a difference taken along a line but an integral taken around a loop --- the $\oint$
that will become this chapter's protagonist, a quantity accumulated all the way around and totted up where it
returns.

Kepler's three laws give the loop its shape, and they are best taken as descriptions of the loop, not
equations to solve. The orbit is an ellipse with the attracting body at a focus,
\begin{equation}
  r = \frac{p}{1 + e\cos\theta};
\end{equation}
equal areas are swept in equal times,
\begin{equation}
  \frac{dA}{dt} = \frac{L}{2m} = \text{const},
\end{equation}
which is the conservation of angular momentum wearing a geometric disguise; and the period and size are
locked together,
\begin{equation}
  T^2 = \frac{4\pi^2}{GM}\,a^3 .
\end{equation}
Each names a feature of the closed loop the world traces and the instrument tallies --- its shape, its rate of
sweep, its size-and-period lock --- a feature we name, not a number we solve for.

Behind the closed shape stands a single function that explains why the loop is a loop. Fold the angular part
of the motion into the radial coordinate alone, and the dynamics is governed by an effective potential,
\begin{equation}
  V_{\text{eff}}(r) = \frac{L^2}{2mr^2} - \frac{GMm}{r},
\end{equation}
two competing terms: the centrifugal barrier $L^2/2mr^2$, which rises steeply as $r \to 0$ and forbids the
body from falling to the center, and the gravitational well $-GMm/r$, which draws it in. Their sum has a
minimum, a basin in $r$, and the orbit is the body rolling in that basin --- resting at the bottom for a
circle, swinging up the slopes for an ellipse, the radius oscillating between a nearest and a farthest turn
while the angle winds around. The whole closed cycle is fixed by two conserved invariants the loop carries, both of them the world's:
the energy $E$, which sets how deep in the basin the motion sits, and the angular momentum $L$, which sets the
height of the centrifugal barrier. Whether the motion is bound or escapes the instrument reads off $E$; the
shape of the conic, off both together.

Two of the instrument's honest floors meet here, and both deserve naming. The closure itself --- the demand
that the loop return to its start --- is a finite count obligation: the instrument verifies directly that the
cycle comes back, a count condition the world enforces and the instrument can check. But the smooth ellipse is
something else, and it is the smooth-shadow ceiling making its first appearance. The continuous curve $r =
p/(1+e\cos\theta)$ is not a thing the instrument holds; it is the smooth shadow that a closed discrete cycle of
counts casts as the cycle is imagined refined --- the shape the loop of counts approaches, ours to draw and
named as a shadow, not the world's to hand over as a fact.
This is the third of the four honest floors, and the rest of the book will rely on it heavily: wherever a
continuum law stands in for a discrete count, the instrument names the bridge an analogy and not a
derivation.

The reading is Kepler's, recast: an orbit is a closed account that must be balanced forever, sustained not by a
force pinning a body in place but by a reconciliation that never finishes --- the body's to keep returning, the
instrument's to keep checking that it has. The loop is the instrument's new home, and the integral around it
the instrument's new tool.

\section{The Angular Momentum Effect: conservation around a loop}
\label{se:angular-momentum}

If the orbit gives the instrument a loop, rotation gives it a tally to carry around one. A rotation conserves
a count, and that count is the angular momentum --- the tally a closed motion carries with it as it goes
around. Written as the cross product of position and momentum,
\begin{equation}
  \mathbf{L} = \mathbf{r} \times \mathbf{p},
\end{equation}
it is the previous chapter's momentum lifted from a line to a loop: where linear momentum balanced across a
boundary, angular momentum is conserved around a cycle. The axis we resolve it about is ours to set; the
conserved tally itself, once an axis is fixed, is the world's. For a rigid body it takes the familiar form
\begin{equation}
  L = I\,\omega, \qquad I \approx M R^2, \qquad \omega = 2\pi f,
\end{equation}
and a hoop of mass $M = 2.0~\mathrm{kg}$ and radius $R = 0.33~\mathrm{m}$ spun at $f = 5~\mathrm{Hz}$ carries
$\omega = 2\pi(5) = 31.4~\mathrm{rad/s}$ and $L = M R^2 \omega \approx 6.8~\mathrm{kg\,m^2/s}$ --- a definite
tally, read off the spin. The conservation law is descriptive bookkeeping,
\begin{equation}
  \frac{dL}{dt} = \tau,
\end{equation}
saying that the tally changes only under a torque and holds fixed when none acts: the loop carries its
angular momentum around unchanged unless something at the boundary turns it. The bookkeeping is ours; that the
tally holds is the world's. The angular coordinate earns its place by compression --- it folds two independent
linear displacements into a single number, carrying half the dimensional burden of the motion it describes.

And the tally is conserved for a reason, not by accident, and naming the reason here gives the book one of its
organizing ideas. Conservation follows from symmetry: this is Noether's theorem, met for the first time.
Because the central well is the same in every direction --- rotate the whole configuration and nothing
physical changes --- the motion carries a rotational symmetry, and to every continuous symmetry Noether's
theorem attaches a conserved count. The count that rotational symmetry conserves is exactly the angular
momentum $\mathbf{L} = \mathbf{r} \times \mathbf{p}$. The instrument does not impose the conservation, nor does
it derive the theorem; it reads the conserved count off a symmetry the world already has. The same principle
returns in the next chapter as the general rule behind every conserved quantity the book meets --- a symmetry
breeds an invariant, and the invariant is a count the instrument can carry around a loop unchanged.

The honest floor is a finite count obligation: a conserved tally the loop carries around, count for count.
The reading is that angular momentum is the conserved count of a rotation --- the rotational echo of the
linear momentum of the previous chapter, balanced now around a cycle rather than across a boundary. The
instrument carries a tally around the loop and finds it unchanged, and the unchanged return is the world's
verdict, not ours: no axis we set, no frame we keep the books in, alters that the tally comes home as it left.
That unchangedness is the conservation --- and it is the calm before the more interesting case, where what the
instrument carries around the loop does \emph{not} come back the same, and the loop returns marked with a
residue the world put there.

\section{The Foucault Effect: precession as holonomy}
\label{se:foucault}

Here is the chapter's payoff, and one of the most important moments in the book: the first time the
instrument carries something around a loop and finds it does not come back the same. Foucault's pendulum is
the demonstration, and the effect deserves to be built slowly, because the capability it introduces ---
holonomy --- is the one every later proof rests on.

Start with the bare idea, on the bench. Give the instrument a tally to carry --- not a number this time but a
\emph{direction}, an arrow it keeps pointed ``the same way'' as it moves. Carry the arrow along a straight
path and back, and it returns pointing exactly as it left: nothing accumulated. Now carry it around a closed
loop that bends --- around a cone, say, or across the curved surface of the Earth --- keeping it as parallel
as the surface allows at every step, never deliberately turning it. When it returns to its starting point, it
is rotated. No one turned it; each step was as straight as the surface permitted; and yet the trip around the
closed loop has left a residue, an angle the arrow has gained for going around. The rule we carry it by, and
the frame we call ``unturned,'' are ours; that it comes back turned all the same is the world's --- the loop's
own curvature, which no rule or frame of ours could have kept straight. That leftover angle is the
\emph{holonomy} --- the residue a closed loop carries, read where the path returns: the angle the world's, the
zero we read it from ours.

The notation names the residue without computing it. Around the loop, the accumulated turning is written
\begin{equation}
  \Delta\phi \;=\; \oint (\text{connection}) \;=\; 2\pi\,(1 - \cos\theta),
\end{equation}
where $\theta$ is the half-angle of the cone the loop traces. This integral is not a number to be ground out;
it is the \emph{name} of what the loop carries --- the holonomy --- and the form $2\pi(1 - \cos\theta)$ is its
\emph{shape}: zero for a flat loop ($\theta = 0$, the arrow returns unturned), growing as the cone sharpens,
exactly the solid angle the loop encloses. The instrument reads the residue off where the arrow returns; the
equation describes that residue, it does not produce it --- ours to write, the world's to have already put
there.

It is worth saying precisely what that connection is, because the same object stands under every proved
residue to come. A connection is a rule for carrying a direction from one point to the next --- the rule of
parallel transport, ``as straight as the surface allows.'' On a flat surface the rule is trivial, and any loop
returns the arrow unturned; on a curved surface it is not, and the failure to return is governed by the
surface's curvature. The holonomy is the loop integral of the connection, and by the theorem of Gauss and
Bonnet that loop integral equals the total curvature the loop encloses,
\begin{equation}
  \oint (\text{connection}) = \iint (\text{curvature}) = \text{the solid angle the loop subtends}.
\end{equation}
For a loop on the unit sphere the enclosed curvature is exactly the area it bounds --- the solid angle --- so
the holonomy is $2\pi(1 - \cos\theta)$, the cap the loop cuts off. The residue the arrow gains is not a
property of the path alone but of the curvature the path encloses, read off the boundary by the connection ---
the connection ours to apply, the curvature the world's to have. This is the precise machinery the proved
chapters will carry: where Foucault's loop encloses a smooth curvature and returns a continuous angle, a
discrete loop will enclose a quantized curvature and return a residue of exactly $\pm 1$.

Foucault's pendulum is this picture made visible. As the Earth turns, the pendulum's swing plane is carried
around a circle of latitude --- a closed loop on a curved surface --- and the plane, kept as parallel as the
motion allows, returns rotated. It precesses, at a rate
\begin{equation}
  \Omega_{\mathrm{prec}} = \Omega \sin\lambda,
\end{equation}
with $\Omega$ the Earth's rotation rate and $\lambda$ the latitude, completing a full turn in one sidereal
day divided by $\sin\lambda$. No torque turns the plane --- there is none --- yet it precesses, because the
loop it is carried around will not let the swing's orientation close. That it precesses is the world's, forced
by the latitude circle it rides; the fixed stars we clock the precession against are the frame we supply. The
precession is the holonomy of that loop, classical and gentle, visible to anyone who watches the pendulum for
an afternoon.

The honest floor is the smooth-shadow analogy the orbit introduced, met now in its sharpest form. The finite
model is a discrete transport-residue --- the failure to close, counted as the swing direction is carried
slot by slot around the loop --- and the continuum formula $2\pi(1 - \cos\theta)$ is the smooth shadow that
residue casts as the slots refine. The bridge between them is an analogy, named and not hidden: the finite
experiment secures a residue, and the precession equation is the continuum that residue points toward, not a
theorem the finite model proves.

And here is the sentence the rest of the book turns on. Foucault is the \emph{gentle, classical first
sighting} of a residue that will return, in the chapters to come, as something the instrument does not
merely shadow but \emph{proves}. The same loop that here leaves a continuous precession will, in an echo
chamber, a diffracting crystal, a rotating ring, a Dirac particle, and a positron's annihilation, carry a
residue of exactly $\pm 1$ --- a sign the instrument can establish as a finite theorem about the loop,
converses and all, not an analogy to a continuum. Foucault shows the reader, in the most familiar setting
there is, what a loop residue \emph{is}; the proofs come later, but the capability is born here. Everything
the instrument will prove about a charged loop, it first learns to see in the slow turning of a pendulum.

\bigskip

With orbits and rotation the instrument has its loop, its tally, and --- decisively --- its holonomy. It can
carry a count around a closed path and read it where it returns: sometimes the count comes back unchanged, and
that is conservation; sometimes it comes back with a residue, and that is holonomy --- and in both the count is
the world's, the frame it is measured against ours. The smooth-shadow ceiling
has debuted, naming the continuum shapes the discrete cycles only approach. Part I closes next with the least
principle and the unseen, where the instrument learns to reconstruct the stretches of motion it did not
directly record --- and to say, honestly, exactly how much of what it reports it has seen and how much it has
inferred.

\bigskip
\begin{center}
\rule{0.4\textwidth}{0.4pt}
\end{center}
\bigskip

\noindent Some accounts do not run in a line from start to finish; they close. The record of such a motion
returns to the very configuration it began in --- a course of events that circles rather than ends, coming
back to its own first entry period after period. That the path closes, and keeps closing, is the world's; the
mark we set at the start, the point we name the beginning of the round, is ours. A closed account is a
different kind of exhibit from a stretch of road between two places: it is a circuit, a thing the record can be
carried around and brought back to where it began. It is the plainest such exhibit Part I enters --- the
account that comes home.
